\ssscp{\tsc{\ili{Central} Flores}}{03-02-03-02}
\citet{Elias2018} shows that there is strong evidence for a node
containing \ili{Rongga}, \ili{Ngadha}, \ili{Nage}, \ili{Keo}, Nga{\Q}o, and \ili{Lio}.
Palu{\Q}e, spoken on a small volcanic island of the same name
to the north of Flores,
is “without a doubt the closest relative of the CF [\tsc{\ili{Central}
Flores}] languages but lies outside the scope of some of the
key innovations defining CF [\tsc{\ili{Central} Flores}] as a subgroup” \citep[85]{Elias2018}.
We label the group containing Palu{\Q}e \tsc{\ili{Central} Flores}
and the node excluding Palu{\Q}e \tsc{Nuclear \ili{Central} Flores}.
The position of \ili{Namut} (a.k.a.
Ngina Manu) with respect to \tsc{West Flores}
and \tsc{\ili{Central} Flores} is unclear due to scarcity of data.
Available data indicates that it shares features
with both groups, though it is perhaps slightly closer to \tsc{\ili{Central} Flores}.

\begin{table}\tcp{\tsc{\ili{Central} Flores} *C > {\0} /{\gap}{\#}}{03-34}
%\begin{adjustbox}{width=\textwidth}
	\begin{threeparttable}\st{2pt}
	\begin{tabular}{lllllll}\lsptoprule
*gloss	&	moon	&	black	&	wind	&	turmeric	&	pith	&	child\\
PMP	&	*bula\tbr{n}	&	*ma-qitə\tbr{m}	&	*haŋi\tbr{n}	&	*kuni\tbr{j}	&	*qunə\tbr{j}\su{†}	&	*ana\tbr{k}\\	\midrule
Palu{\Q}e	&	{\hr}vula	&	\cg{(more)}	&	{\hh}aŋi	&		&	{\hq}une	&	{\hr}ana\\
\ili{Lio}	&	{\hr}ʋula	&	{\hr}mite	&	{\hh}aŋi	&	{\hr}kune	&	{\hq}une	&	{\hr}ana\\
\ili{Ende}	&	{\hr}wuɹa	&	{\hr}mite	&	{\hh}aŋi	&	{\hr}kune	&	{\hr}ʔune	&	{\hr}ana\\
\ili{Central} \ili{Nage}	&	{\hr}wula	&	{\hr}mite	&	\hp{h}\cg{(waa)}	&		&		&	{\hr}ana\\
Nga{\Q}o, \ili{Oja}	&	{\hr}ʋud͡ʎ̝a	&	{\hr}mite	&	\hp{h}\cg{(ʋiyo)}	&		&		&	{\hr}ana\\
Nga{\Q}o, Watumite	&	{\hr}ʋud͡ʎ̝a	&	{\hr}mite	&	\hp{h}\cg{(ʋiɹo)}	&		&		&	{\hr}ana\\
East \ili{Nage}	&	{\hr}ʋuda	&	{\hr}mite	&	{\hh}aŋi	&		&		&	{\hr}ana\\
\ili{Keo}, Udiworowatu	&	{\hr}wuda	&	{\hr}mite	&	{\hh}aŋi	&	{\hr}kune	&	{\hr}ʔone	&	{\hr}ana\\
\ili{Ngadha}	&	{\hr}vula	&	{\hr}mite	&	{\hh}aŋi\su{‡}	&	{\hr}kune	&	{\hq}une	&	{\hr}ana\\
\ili{Rongga}	&	{\hr}wula	&	{\hr}mite	&	{\hh}aŋi	&	{\hr}kune	&		&	{\hr}ana\\	\midrule
\ili{Komodo}	&	{\hr}wulaŋ	&	\cg{(mureŋ)}	&	{\hh}aŋiŋ	&	{\hr}gunis	&		&	{\hr}anaʔ\\
\ili{Manggarai}	&	{\hr}ʋulaŋ	&	{\hr}mitəŋ	&	\hp{h}\cg{(buru)}	&	{\hr}ʋunis	&	{\hq}unər	&	{\hr}anak\\
		\lspbottomrule
	\end{tabular}
		\begin{tablenotes}\tnsize
			\item[†]	{All \tsc{\ili{Central} Flores} reflexes
								with the exception of \ili{Ngadha}
								are glossed as ‘inside’.}
			\item[‡]	{\ili{Ngadha} \it{aŋi} is glossed as ‘weak wind’ \citep[199]{Arndt1961}.}
		\end{tablenotes}
	\end{threeparttable}
%\end{adjustbox}
\end{table}

The main sound change which unifies \tsc{\ili{Central} Flores}
and sets it apart from \tsc{West Flores} is the loss of word
final consonants.\footnote{
		Palu{\Q}e has some instances of word-final \it{n}
		which may the third person genitive enclitic \it{=n} in many cases.
		However, even Palu{\Q}e has lost root-final *n in
		reflexes of *bulan > \it{vula} ‘moon’,
		*haŋin > \it{aŋi} ‘wind’,
		*quzan > \it{ura} ‘rain’,
		*ŋajan > \it{ŋara} ‘name’,
		and *zalan > \it{lala} ‘path’.}
Examples of final *C > {\0}
/{\gap}{\#} where a final consonant is retained in \tsc{West
Flores} are given in \trf{03-34}.
Additional examples occur in the reflexes of *(ŋ)ijuŋ ‘nose’
and *ŋajan ‘name’ in \trf{03-26} (\prf{tab:03-26}),
as well as *quzan ‘rain’
and *zalan ‘path’ in \trf{03-27} (\prf{tab:03-27}).

An additional change differentiating \tsc{\ili{Central} Flores}
from \tsc{West Flores} is final *-aw > \it{a} in at l\ili{east} four
words where \tsc{West Flores} has *-aw > \it{o}.
These words are given in \trf{03-35}.
However, *-aw > \it{o} also occurs in \tsc{\ili{Central} Flores}.
Two examples are *panaw > \ili{Lio},
\ili{Ende}, Nga{\Q}o, East \ili{Nage}, \ili{Ngadha} \it{pano} ‘Tinea flava’,
as well as reflexes of *babaw ‘upper’ given in \trf{03-01} (\prf{tab:03-01}).

\begin{table}
	\centering\tcp{\tsc{\ili{Central} Flores} *-aw > \it{a}}{03-35}
	\st{6pt}\begin{tabular}{lllll}\lsptoprule
*gloss	&	sun, day	&	go	&	go	&	steal\\
PMP	&	*qaləj\tbr{aw}	&	*pan\tbr{aw}	&	*lak\tbr{aw}	&	*nak\tbr{aw}\\	\midrule
Palu{\Q}e	&	{\hqa}ʔər\tbr{a}	&	{\hr}pan\tbr{a}	&		&	{\hr}nak\tbr{a}\\
\ili{Lio}	&	{\hqa}ləʤ\tbr{a}	&	{\hr}mban\tbr{a}	&		&	{\hr}nak\tbr{a}\\
\ili{Ende}	&	{\hqa}ɹər\tbr{a}	&	{\hr}mban\tbr{a}	&		&	{\hr}nak\tbr{a}\\
\ili{Central} \ili{Nage}	&	{\hqa}ləz\tbr{a}	&	{\hr}mban\tbr{a}	&		&	{\hr}nak\tbr{a}\\
Nga{\Q}o, \ili{Oja}	&	{\hqa}d͡ʎ̝əy\tbr{a}	&	{\hr}mban\tbr{a}	&		&	{\hr}nak\tbr{a}\\
Nga{\Q}o, Watumite	&	{\hqa}d͡ʎ̝əɹ\tbr{a}	&	{\hr}pan\tbr{a}	&		&	{\hr}nak\tbr{a}\\
East \ili{Nage}	&	{\hqa}dəɹ\tbr{a}	&	{\hr}ban\tbr{a}	&		&	{\hr}nak\tbr{a}\\
\ili{Keo}, Udiworowatu	&	{\hqa}dər\tbr{a}	&	{\hr}mban\tbr{a}	&		&	{\hr}nak\tbr{a}\\
\ili{Ngadha}	&	{\hqa}ləʣ\tbr{a}	&		&	{\hr}laʔ\tbr{a}	&	{\hr}nak\tbr{a}\\
\ili{Rongga}	&	{\hqa}ləɹ\tbr{a}	&		&	{\hr}laʔ\tbr{a}	&	{\hr}nak\tbr{a}\\	\midrule
\ili{Waerana}	&	{\hqa}ləzoʔ	&	\cg{(lətan)}	&	{\hr}lakoʔ	&	{\hr}takar\\
\ili{Manggarai}	&	{\hqa}ləso	&	\cg{(ŋoo)}	&	{\hr}lako	&	{\hr}tako\\
		\lspbottomrule
	\end{tabular}
\end{table}

\clearpage
All \tsc{\ili{Central} Flores} languages except Palu{\Q}e can be placed
in a \tsc{Nuclear \ili{Central} Flores} node.
The evidence that \citet[119--147]{Elias2018} amasses for this
node includes shared reconstructable mixed quinary-decimal numerals
(\trf{03-36}), shared reconstructable paradigm of directional terms (\trf{03-37}),
lexically specific irregular sound changes (\trf{03-38}),
and lexical replacements in basic vocabulary (\trf{03-39}).
Most of the innovative numerals are also present in \ili{Waerana}.
This could be due to contact between \ili{Waerana}
and \tsc{\ili{Central} Flores} languages,
or \ili{Waerana} may be transitional between \tsc{West Flores}
and \tsc{\ili{Central} Flores}.

\begin{table}
%\tcp{\tsc{Nuclear \ili{Central} Flores} innovative numerals \citep[121]{Elias2018}}{03-36}
\tcp{\tsc{Nuclear \ili{Central} Flores} innovative numerals}{03-36}
%\begin{adjustbox}{width=\textwidth}
	\begin{threeparttable}\st{2pt}
	\begin{tabular}{llllll}\lsptoprule
gloss	&	four	&	six	&	seven	&	eight	&	nine\\
Proto-NCF	&	**(sa=)wutu\su{†}	&	**lima əsa	&	**lima dua	&	**dua mbutu\su{†}	&	**təra əsa\\ \midrule
\ili{Lio}	&	{\hr\hr}sutu	&	{\hr\hr}lima əsa	&	{\hr\hr}lima rua	&	{\hr\hr}rua mbutu	&	{\hr\hr}təra əsa\\
\ili{Ende}	&	{\hr\hr}wutu	&	{\hr\hr}ɹima ʔəsa	&	{\hr\hr}ɹima rua	&	{\hr\hr}rua mbutu	&	{\hr\hr}təra ʔəsa\\
C. \ili{Nage}	&	{\hr\hr}wutu	&	{\hr\hr}lima əsa	&	{\hr\hr}lima zua	&	{\hr\hr}ɹua butu	&	{\hr\hr}taa əsa\\
Nga{\Q}o, \ili{Oja}	&	{\hr\hr}ʋutu	&	{\hr\hr}d͡ʎ̝ima əsa	&	{\hr\hr}d͡ʎ̝ima rua	&	{\hr\hr}rua butu	&	{\hr\hr}təra əsa\\
Nga{\Q}o, Wat.	&	{\hr\hr}ʋutu	&	{\hr\hr}d͡ʎ̝ima əsa	&	{\hr\hr}d͡ʎ̝ima ɹua	&	{\hr\hr}ɹua butu	&	{\hr\hr}təra əsa\\
East \ili{Nage}	&	{\hr\hr}ʋutu	&	{\hr\hr}dima əsa	&	{\hr\hr}dima ɹua	&	{\hr\hr}ɹua mbutu	&	{\hr\hr}təra əsa\\
\ili{Keo}, Ud.	&	{\hr\hr}wutu	&	{\hr\hr}dima ʔəsa	&	{\hr\hr}dima rua	&	{\hr\hr}rua mbutu	&	{\hr\hr}təra ʔəsa\\
\ili{Ngadha}	&	{\hr\hr}vutu	&	{\hr\hr}lima əsa	&	{\hr\hr}lima zua	&	{\hr\hr}rua butu	&	{\hr\hr}təra əsa\\
\ili{Rongga}	&	{\hr\hr}wutu	&	{\hr\hr}lima əsa	&	{\hr\hr}lima ɹua	&	{\hr\hr}ɹua mbutu	&	{\hr\hr}tara əsa\\ \hhline{~~~~~-}
%\ili{Waerana}	&	{\hr}sutuʔ	{\cgr}&	{\hr}lima saʔ	{\cgr}&	{\hr}lima zuaʔ	{\cgr}&	{\hr}ara butuʔ	{\cgr}&	{\hr}siok\\
\ili{Waerana}	&	{\hr\hr}sutuʔ	&	{\hr\hr}lima saʔ	&	{\hr\hr}lima zuaʔ	&	\mc{1}{l|}{{\hr\hr}ara butuʔ}	&	{\hr\hr}siok\\ \midrule
PMP	&	\hr*əpat	&	\hr*ənəm	&	\hr*pitu	&	\hr*walu	&	\hr*siwa\\ \midrule
Palu{\Q}e	&	{\hr\hr}ɓaa	&	{\hr\hr}ʔəne	&	{\hr\hr}ɓitu	&	{\hr\hr}valu	&	{\hr\hr}iwa\\
%\ili{Rembong}	&	{\hr\hr}paat	&	{\hr\hr}noon	&	{\hr\hr}pituʔ	&	{\hr\hr}waluʔ	&	{\hr\hr}siwaʔ\\
\ili{Manggarai}	&	{\hr\hr}paat	&	{\hr\hr}ənəm	&	{\hr\hr}pitu	&	{\hr\hr}alo	&	{\hr\hr}ʧiok\\
		\lspbottomrule
	\end{tabular}
		\begin{tablenotes}\tnsize
			\item[†]	{
								\tsc{Nuclear \ili{Central} Flores} **wutu ‘four’
								and the second part of **dua mbutu ‘eight’ are
								connected with PCEMP **butu ‘group,
								crowd, flock, school, bunch,
								cluster’ \citep{BlustTrussel2020} and/or Proto-\ili{Alor}-\ili{Pantar} (non-Austronesian) *buta ‘four’ \citep[283]{SchapperKlamer2017}.
								\tsc{Flores-Lembata} languages (except \ili{Sika})
								have a reflex of **butu in their numeral for ‘eight’; \ili{Lamaholot}
								\it{buto} ‘eight’ and \ili{Kedang} \it{butu rai} ‘eight’ \citep[366f]{Fricke2019}.
								\ili{Sika} \it{hutu} ‘four’ may be borrowed from \ili{Lio} \it{sutu}
								with regular \ili{Sika} *s > \it{h}.}
		\end{tablenotes}
	\end{threeparttable}
%\end{adjustbox}
\end{table}

\begin{table}
%\caption[\tsc{Nuclear \ili{Central} Flores} paradigm of directional terms (Elias 2018:122)]
\caption[\tsc{Nuclear \ili{Central} Flores} paradigm of directional terms]
{\tsc{Nuclear \ili{Central} Flores} paradigm of directional terms\su{†}}
\label{tab:03-37}
\begin{adjustbox}{width=\textwidth}
	\begin{threeparttable}\st{2pt}
	\begin{tabular}{lllllll}\lsptoprule
PMP	&		&		&	*lahud	&	*i atas	&	*i daləm	&	\\
NCF gloss	&	{\hr}down	&	{\hr}up	&	{\hr}seawards	&	{\hr}inwards, above	&	{\hr}below, west	&	{\hr}\ili{east}\su{‖}\\
Proto-NCF	&	**ʤili\su{‡}	&	**ʤele\su{‡}	&	**lau	&	**ʤeta\su{‡}	&	**ʤale\su{‡}	&	**məna\\	\midrule
\ili{Lio}	&	{\hr}\cg{(ɣawa)}	&	\hp{**}ɣele	&	\hp{**}lau	&	\hp{**}ɣeta	&	\hp{**}ɣale	&	\hp{**}məna\\
\ili{Ende}	&	\hp{**}riɹi	&	\hp{**}reɹe	&	\hp{**}ɹau	&	\hp{**}reta	&	\hp{**}raɹe	&	\hp{**}məna\\
C. \ili{Nage}	&	\hp{**}zili	&	\hp{**}zele	&	\hp{**}lau	&	\hp{**}zeta	&	\hp{**}zale	&	\hp{**}məna\\
East \ili{Nage}	&	\hp{**}ɹidi	&		&		&	\hp{**}ɹeta	&	\hp{**}ɹade	&	\\
\ili{Keo}, Ud.	&	\hp{**}ridi	&	\hp{**}rede	&	\hp{**}dau	&	\hp{**}reta	&	\hp{**}rade	&	\hp{**}məna\\
\ili{Ngadha}	&	\hp{**}zili	&	\hp{**}zele	&	\hp{**}lau	&	\hp{**}zeta	&	\hp{**}zale	&	\hp{**}məna\\
\ili{Rongga}	&	\hp{**}ɹili	&	\hp{**}ɹele	&	\hp{**}lau	&	\hp{**}ɹeta	&	\hp{**}ɹale	&	\hp{**}məna\\
Palu{\Q}e	&		&		&		&	\hp{**}reta	&	{\hr}\cg{(lae)}	&	\\	\midrule
\ili{Waerana}	&		&		&		&	\hp{**}etaʔ	&	{\hr}\cg{(waaʔ)}, aleʔ\su{Δ}	&	{\hr}\cg{(awo)}\\
\ili{Manggarai}	&	{\hr}\cg{(ʋaa)}	&	{\hr}\cg{(eta)}	&	\hp{**}lau	&	\hp{**}eta	&	{\hr}\cg{(vaa)}, sale\su{Δ}	&	{\hr}\cg{(avo)}\\
		\lspbottomrule
	\end{tabular}
		\begin{tablenotes}\tnsize
			\item[†]	{Although four of the terms in this paradigm are inheritances
								PMP and/or more widely distributed
								their use in a paradigm is seen as
								innovative by \citet{Elias2018}.}
			\item[‡]	{The initial consonants of **ʤili,
								**ʤele, **ʤeta, and **ʤale is the same as the result of the merger of *z/*j.
								We follow \citet{Elias2018} in representing this consonant as **ʤ.}
			\item[Δ]	{\ili{Waerana} \it{aleʔ}
								and \ili{Manggarai} \it{sale} are ‘west’,
								while \it{waaʔ} and \it{vaa} are ‘below’.}
			%\item[Δ]	{\ili{Manggarai} \it{sale} is ‘west’,
								%while \it{waaʔ} and \it{vaa} are ‘below’.}
			\item[‖]	{Cognate forms of **məna meaning `front'
								occur in many other languages of Wallacea.
								For example: \ili{Buru} \it{mena-n} ‘front’ and
								\ili{Alune} \it{kwali mena} ‘older same sex sibling’
								(lit. same sex sibling in front),
								\ili{Yalahatan}, \it{mina} ‘front, first, eldest’.
								\ili{Larike}: \it{mina} ‘in front.
								Compare \citeauthor{Stresemann1927}'s
								(\citeyear[113]{Stresemann1927})
								Proto-\ili{Central} Maluku **ma-ina ‘in front of’.
								%**ma-ina (p.113 pdfp. 62)
								Several directional systems in the
								region are anchored around
								the speaker facing \ili{east}.
								So, front = \ili{east}, left = north,
								right = south, back = west.}
		\end{tablenotes}
	\end{threeparttable}
\end{adjustbox}
\end{table}

\begin{table}\tcp{\tsc{Nuclear \ili{Central} Flores} lexically specific sound changes}{03-38}
\begin{adjustbox}{width=\textwidth}
	\begin{threeparttable}\st{2pt}
	\begin{tabular}{lllllll}\lsptoprule
local gloss	&	flower	&	foam	&	tail	&	bad, evil	&	wife	&	head hair\\
PMP	&	*b\tbr{u}ŋa	&	*b\tbr{u}jəq	&	*\tbr{i}kuR	&	*z\tbr{a}q\tbr{a}t	&	*\tbr{b}ahi	&	*\tbr{b}uhək\\
change	&	*u > **o	&	*u > **o	&	*i > **e	&	*a > **e	&	*b > **f\su{†}	&	*b > **f\su{†}\\
Proto-NCF	&	**w\tbr{o}ŋa\su{†}	&	**w\tbr{o}da	&	**\tbr{e}ko	&	**z\tbr{e}ʔ\tbr{e}	&	**\tbr{f}ai	&	**\tbr{f}uu\\	\midrule
\ili{Lio}	&	\hp{**}ʋ\tbr{o}ŋa	&	\hp{**}ʋ\tbr{o}ra	&	\hp{**}\tbr{e}ko	&	\hp{**}r\tbr{e}ʔ\tbr{e}	&	\hp{**}\tbr{f}ai	&	\hp{**}\tbr{f}uu\\
\ili{Ende}	&	\hp{**}w\tbr{o}ŋa	&	\hp{**}w\tbr{o}ra	&	\hp{**}ʔ\tbr{e}ko	&	\hp{**}r\tbr{e}ʔ\tbr{e}	&	\hp{**}\tbr{h}ai	&	\hp{**}\tbr{h}uu\\
C. \ili{Nage}	&	\hp{**}w\tbr{o}ŋa	&		&	\hp{**}\tbr{e}ko	&	\hp{**}\tbr{e}ʔ\tbr{e}\su{Δ}	&	\hp{**}\tbr{f}ai	&	\hp{**}\tbr{f}uu\\
Nga{\Q}o, \ili{Oja}	&	\hp{**}ʋ\tbr{o}ŋa	&	\hp{**}ʋ\tbr{o}ya	&	\hp{**}\tbr{e}ko	&	\hp{**}y\tbr{e}ʔ\tbr{e}	&	\hp{**}\tbr{f}ai	&	\hp{**}\tbr{f}uu\\
Nga{\Q}o, Wat.	&	\hp{**}ʋ\tbr{o}ŋa	&	\hp{**}ʋ\tbr{o}ɹa	&	\hp{**}\tbr{e}ko	&	\hp{**}y\tbr{e}ʔ\tbr{e}\su{Δ}	&	\hp{**}\tbr{f}ai	&	\hp{**}\tbr{f}uu\\
East \ili{Nage}	&	\hp{**}ʋ\tbr{o}ŋa	&	\hp{**}ʋ\tbr{o}ɹa	&	\hp{**}\tbr{e}ko	&	\hp{**}l\tbr{e}ʔ\tbr{e}\su{Δ}	&	\hp{**}\tbr{f}ai	&	\hp{**}\tbr{f}uu\\
\ili{Keo}, Ud.	&	\hp{**}w\tbr{o}ŋa	&		&	\hp{**}ʔ\tbr{e}ko	&	\hp{**}r\tbr{e}ʔ\tbr{e}	&	\hp{**}\tbr{f}ai	&	\hp{**}\tbr{f}uu\\
\ili{Ngadha}	&	\hp{**}v\tbr{o}ŋa	&	\hp{**}v\tbr{o}za\su{‡}	&	\hp{**}\tbr{e}ko	&	\hp{**}z\tbr{e}ʔ\tbr{e}	&	\hp{**}\tbr{f}ai	&	\hp{**}\tbr{f}uu\\
\ili{Rongga}	&	{\hr}\cg{(lando)}	&	\hp{**}w\tbr{o}ɹa	&	\hp{**}\tbr{e}ko	&	\hp{**}ɹ\tbr{e}ʔ\tbr{e}	&	\hp{**}\tbr{f}ai	&	\hp{**}\tbr{f}uu\\	\midrule
Palu{\Q}e	&	{\hr}\cg{(ʧoa)}	&	{\hr}\cg{(koa)}	&	\hp{**}iʔo	&	{\hr}\cg{(kaʔa mbola)}	&	\hp{**}vai	&	{\hr}\cg{(lolo)}\\
\ili{Waerana}	&	\hp{**}buŋa	&	\hp{**}wuzaʔ	&	\hp{**}ikoʔ	&	\hp{**}daʔat	&	\hp{**}wai	&	\hp{**}wuuk\\
%\ili{Rembong}	&	\hp{**}wuŋa, buŋa	&	\hp{**}busaʔ	&	\hp{**}iko, ikoʔ	&	\hp{**}daʔat	&	\hp{**}wai	&	\hp{**}wuuk, wouk\\
\ili{Manggarai}	&	{\hr}\cg{(ʋela)}	&	\hp{**}ʋusa	&	\hp{**}iko	&	\hp{**}daʔat	&	{\hr}\cg{(ʋina)}	&	\hp{**}ʋuuk\\
		\lspbottomrule
	\end{tabular}
		\begin{tablenotes}\tnsize
			\item[†]	{Irregular *b > **f probably went through intermediate **v,
								retained in Palu{\Q}e \it{vai} ‘wife’.}
			\item[‡]	{\ili{Ngadha} also has \it{fusa} ‘foam’ which may be from the PMP doublet *busa.}
			\item[Δ]	{\ili{Central} \ili{Nage},
								Watumite Nga{\Q}o, and East \ili{Nage} have irregular initial consonant reflexes.}
		\end{tablenotes}
	\end{threeparttable}
\end{adjustbox}
\end{table}

\begin{table}
%\caption[\tsc{Nuclear \ili{Central} Flores} lexical replacement (Elias 2018:124)]
\caption[\tsc{Nuclear \ili{Central} Flores} lexical replacement]
{\tsc{Nuclear \ili{Central} Flores} lexical replacement\su{†}}
\label{tab:03-39}
%\begin{adjustbox}{width=\textwidth}
	\begin{threeparttable}\st{2pt}
	\begin{tabular}{llllllll}\lsptoprule
local gloss	&	night	&	know	&	heavy	&	betel nut	&	dig	&	back	&	near\\
Proto-NCF	&	**kobe	&	**mbeʔo	&	**ndate	&	**kleu	&	**koe	&	**loŋgo	&	**weʔe\\	\midrule
\ili{Lio}	&	\hp{**}kobe	&	\hp{**}mbeʔo	&	\hp{**}ndate	&	\hp{**}keu	&	\hp{**}koe	&	\hp{**}loŋgo	&	\hp{**}ʋeʔe\\
\ili{Ende}	&	\hp{**}kombe	&	\hp{**}mbeʔo	&	\hp{**}ndate	&	\hp{**}eu	&	\hp{**}koe	&	\hp{**}ɹoŋgo	&	\hp{**}weʔe\\
C. \ili{Nage}	&	\hp{**}kobe	&	\hp{**}beʔo	&		&	\hp{**}heu	&	\hp{**}koe	&	\hp{**}logo	&	\hp{**}weʔe\\
Nga{\Q}o, \ili{Oja}	&	\hp{**}kombe	&	\hp{**}mbeʔo	&	\hp{**}ndate	&	\hp{**}yeu	&	\hp{**}koe	&	\hp{**}d͡ʎ̝oŋgo	&	\hp{**}ʋeʔe\\
Nga{\Q}o, Wat.	&	\hp{**}kob̥e\su{‡}	&	\hp{**}b̥eʔo\su{‡}	&	\hp{**}ndate	&	\hp{**}yeu	&	\hp{**}koe	&	\hp{**}d͡ʎ̝og̊o\su{‡}	&	\hp{**}ʋeʔe\\
East \ili{Nage}	&	\hp{**}kombe	&	\hp{**}mbeʔo	&	\hp{**}ndate	&	\hp{**}leu	&	\hp{**}koe	&	\hp{**}doŋgo	&	\hp{**}ʋeʔe\\
\ili{Keo}, Ud.	&	\hp{**}kombe	&	\hp{**}mbeʔo	&	\hp{**}ndate	&	\hp{**}heu	&	\hp{**}koe	&	\hp{**}doŋgo	&	\hp{**}weʔe\\
\ili{Ngadha}	&	\hp{**}kobe	&	\hp{**}beʔo	&	\hp{**}date	&	\hp{**}heu	&	\hp{**}koe	&	\hp{**}logo	&	\hp{**}veʔe\\
\ili{Rongga}	&	\hp{**}kombe	&	\hp{**}mbeʔo	&	\hp{**}ndate	&	\hp{**}heu	&	\hp{**}koe	&	\hp{**}loŋgo	&	\hp{**}weʔe\\	\midrule\midrule
PMP	&	*bəRŋi	&	*taqu	&	*bəRəqat	&	*buaq	&	*kali	&	*likud	&	*adani\\ \midrule
Palu{\Q}e	&	\cg{(məre)}	&	\cg{(ʧuʔu)}	&	{\hr}pəʤa	&	{\hr}vua	&	{\hr}kali	&	\cg{(tola)}	&	{\hr}təni\\
\ili{Rembong}	&	\cg{(beweʔ)}	&	\cg{(lətoʔ)}	&	{\hr}bərat	&	\cg{(rasi)}	&	\cg{(kakek)}	&	\cg{(muzi)}	&	\cg{(soʔek)}\\
\ili{Waerana}	&	{\hr}bəŋuʔ	&	\cg{(soʔuk)}	&	\cg{(məndoʔ)}	&	\cg{(basi)}	&	{\hr}ɣaliʔ	&	\cg{(atur)}	&	\cg{(soʔek)}\\
Kepo{\Q}	&	\cg{(mane)}	&	\cg{(tiŋon)}	&	\cg{(məndo)}	&	\cg{(kabe)}	&	\cg{(sake)}	&		&	\cg{(ruis)}\\
\ili{Manggarai}	&	\cg{(ʋie)}	&	\cg{(pəʧiŋ)}	&	\cg{(məndo)}	&	\cg{(raʧi)}	&	\cg{(ʧake)}	&	\cg{(toni)}	&	\cg{(baliŋ)}\\
		\lspbottomrule
	\end{tabular}
		\begin{tablenotes}\tnsize
			\item[†]	{We give only seven of the eleven lexical replacements in \citet{Elias2018}.
								The other five putative lexical replacements have cognates elsewhere
								or a problematic in other respects.
								**toro ‘red’ (\ili{Waerana} \it{toroŋ} ‘red’,
								\ili{Kodi} \it{ɗoro} ‘red’) **liru ‘sky’ (\ili{Havu},
								\ili{Dhao} \it{liru} ‘sky’, possibly \ili{Bima} \it{liro} ‘sun’), **pawe ‘good’
								in place of *diqa (\ili{Rongga} retains \it{ɹiʔa},
								and \ili{Keo} retains \it{ria} alongside innovative \it{pawe}),
								and **teʔe in place of *təpiR for ‘mat’ (possible irregular reflex).
								}
			\item[‡]	{Watumite Nga{\Q}o has contrast between aspirated /pʰ/,
								voiceless /p/, voiced /b/,
								and imploded /ɓ/, as well as between /kʰ/, /k/, and /g/
								\citep[78]{Elias2018}.
								We follow \citeauthor{Elias2018} in using the symbols ‹b̥›
								and ‹g̊› for the voiceless unaspirated plosives [p] and [k] respectively,
								and ‹p› and ‹k› for aspirated plosives [pʰ] and [kʰ].
								In other languages of the region aspiration is
								a non-contrastive feature of voiceless plosives.}
		\end{tablenotes}
	\end{threeparttable}
%\end{adjustbox}
\end{table}

\clearpage
\citet{Elias2018} argues convincingly that while \tsc{Nuclear
\ili{Central} Flores} is a well-defined subgroup,
it is difficult to propose lower level mutually exclusive subgroups
due to the non-tree like pattern of overlapping innovations.
Nonetheless, at l\ili{east} one further tree-like split can be identified
within \tsc{Nuclear \ili{Central} Flores}.
All languages except \ili{Lio} usually merge **z (from *z/*j) with *d > **z.
\ili{Lio}, on the other hand has distinct reflexes with *d > \it{r}
and **z > \it{ʤ}.
Examples of this latter sound change are given in \trf{03-40}.

\begin{table}\tcp{\tsc{\ili{Central} East Flores} **z/*d > **z}{03-40}
%\begin{adjustbox}{width=\textwidth}
	\begin{threeparttable}\st{5.5pt}
	\begin{tabular}{lllllll}\lsptoprule
local gloss	&	path	&	sun, day	&	name	&	new	&	bathe	&	live\\
PMP	&	*\tbr{z}alan	&	*qalə\tbr{j}aw	&	*ŋa\tbr{j}an	&	*ma-u\tbr{d}əhi	&	*\tbr{d}iRus	&	*ma-qu\tbr{d}ip\\ \midrule
Palu{\Q}e	&	{\hr}\tbr{l}ala	&	{\hr}ʔə\tbr{r}a	&	{\hr}ŋa\tbr{r}a	&	{\hr}mu\tbr{r}i	&	{\hr}\tbr{t}io	&	\\
\ili{Lio}	&	{\hr}\tbr{ʤ}ala	&	{\hr}lə\tbr{ʤ}a	&	{\hr}ŋa\tbr{ʤ}a	&	{\hr}mu\tbr{r}i	&	{\hr}\tbr{r}io	&	{\hr}mu\tbr{r}i\\ \midrule
\ili{Ende}	&	{\hr}\tbr{r}aɹa	&	{\hr}ɹə\tbr{r}a	&	{\hr}ŋa\tbr{r}a	&	{\hr}mu\tbr{r}i	&	{\hr}\tbr{r}io	&	{\hr}mu\tbr{r}i\\
C. \ili{Nage}	&	{\hr}\tbr{z}ala	&	{\hr}lə\tbr{z}a	&	{\hr}ŋa\tbr{z}a	&	{\hr}mu\tbr{z}i	&	{\hr}\tbr{z}io	&	\\
Nga{\Q}o, \ili{Oja}	&	{\hr}\tbr{y}ad͡ʎ̝a	&	{\hr}d͡ʎ̝ə\tbr{y}a	&	{\hr}ŋa\tbr{y}a	&	{\hr}mu\tbr{y}i	&	{\hr}\tbr{y}io	&	{\hr}mu\tbr{y}i\\
Nga{\Q}o, Wat.	&	{\hr}\tbr{ɹ}ad͡ʎ̝a	&	{\hr}d͡ʎ̝ə\tbr{ɹ}a	&	{\hr}ŋa\tbr{ɹ}a	&	{\hr}mu\tbr{ɹ}i	&	{\hr}\tbr{ɹ}io	&	{\hr}mu\tbr{ɹ}i\\
East \ili{Nage}	&	{\hr}\tbr{ɹ}ada	&	{\hr}də\tbr{ɹ}a	&	{\hr}ŋa\tbr{ɹ}a	&	{\hr}mu\tbr{ɹ}i	&	{\hr}\tbr{ɹ}io	&	{\hr}mu\tbr{ɹ}i\\
\ili{Keo}, Ud.	&	{\hr}\tbr{r}ada	&	{\hr}də\tbr{r}a	&	{\hr}ŋa\tbr{r}a	&	{\hr}mu\tbr{r}i	&	{\hr}\tbr{r}io	&	{\hr}mu\tbr{r}i\\
\ili{Ngadha}	&	{\hr}\tbr{ʣ}ala	&	{\hr}lə\tbr{ʣ}a	&	{\hr}ŋa\tbr{ʣ}a	&	{\hr}mu\tbr{ʣ}i	&	{\hr}\tbr{ʣ}io	&	{\hr}mu\tbr{ʣ}i\\
\ili{Rongga}	&	{\hr}la\tbr{ɹ}a\su{†}	&	{\hr}lə\tbr{ɹ}a	&	{\hr}ŋa\tbr{ɹ}a	&	{\hr}mu\tbr{ɹ}i	&	{\hr}\tbr{ɹ}io	&	{\hr}mu\tbr{ɹ}i\\
		\lspbottomrule
	\end{tabular}
		\begin{tablenotes}\tnsize
			\item[†]	{\ili{Rongga} \it{laɹa} ‘path’ has consonant metathesis.}
		\end{tablenotes}
	\end{threeparttable}
%\end{adjustbox}
\end{table}

